% Options for packages loaded elsewhere
\PassOptionsToPackage{unicode}{hyperref}
\PassOptionsToPackage{hyphens}{url}
\PassOptionsToPackage{dvipsnames,svgnames,x11names}{xcolor}
\documentclass[
  spanish,
  10pt,
  a4paper,
]{article}
\usepackage{xcolor}
\usepackage[top=2.2cm,bottom=2.2cm,left=2.1cm,right=2.1cm]{geometry}
\usepackage{amsmath,amssymb}
\setcounter{secnumdepth}{5}
\usepackage{iftex}
\ifPDFTeX
  \usepackage[T1]{fontenc}
  \usepackage[utf8]{inputenc}
  \usepackage{textcomp} % provide euro and other symbols
\else % if luatex or xetex
  \usepackage{unicode-math} % this also loads fontspec
  \defaultfontfeatures{Scale=MatchLowercase}
  \defaultfontfeatures[\rmfamily]{Ligatures=TeX,Scale=1}
\fi
\usepackage{lmodern}
\ifPDFTeX\else
  % xetex/luatex font selection
\fi
% Use upquote if available, for straight quotes in verbatim environments
\IfFileExists{upquote.sty}{\usepackage{upquote}}{}
\IfFileExists{microtype.sty}{% use microtype if available
  \usepackage[]{microtype}
  \UseMicrotypeSet[protrusion]{basicmath} % disable protrusion for tt fonts
}{}
\makeatletter
\@ifundefined{KOMAClassName}{% if non-KOMA class
  \IfFileExists{parskip.sty}{%
    \usepackage{parskip}
  }{% else
    \setlength{\parindent}{0pt}
    \setlength{\parskip}{6pt plus 2pt minus 1pt}}
}{% if KOMA class
  \KOMAoptions{parskip=half}}
\makeatother
\usepackage{color}
\usepackage{fancyvrb}
\newcommand{\VerbBar}{|}
\newcommand{\VERB}{\Verb[commandchars=\\\{\}]}
\DefineVerbatimEnvironment{Highlighting}{Verbatim}{commandchars=\\\{\}}
% Add ',fontsize=\small' for more characters per line
\newenvironment{Shaded}{}{}
\newcommand{\AlertTok}[1]{\textcolor[rgb]{1.00,0.00,0.00}{\textbf{#1}}}
\newcommand{\AnnotationTok}[1]{\textcolor[rgb]{0.38,0.63,0.69}{\textbf{\textit{#1}}}}
\newcommand{\AttributeTok}[1]{\textcolor[rgb]{0.49,0.56,0.16}{#1}}
\newcommand{\BaseNTok}[1]{\textcolor[rgb]{0.25,0.63,0.44}{#1}}
\newcommand{\BuiltInTok}[1]{\textcolor[rgb]{0.00,0.50,0.00}{#1}}
\newcommand{\CharTok}[1]{\textcolor[rgb]{0.25,0.44,0.63}{#1}}
\newcommand{\CommentTok}[1]{\textcolor[rgb]{0.38,0.63,0.69}{\textit{#1}}}
\newcommand{\CommentVarTok}[1]{\textcolor[rgb]{0.38,0.63,0.69}{\textbf{\textit{#1}}}}
\newcommand{\ConstantTok}[1]{\textcolor[rgb]{0.53,0.00,0.00}{#1}}
\newcommand{\ControlFlowTok}[1]{\textcolor[rgb]{0.00,0.44,0.13}{\textbf{#1}}}
\newcommand{\DataTypeTok}[1]{\textcolor[rgb]{0.56,0.13,0.00}{#1}}
\newcommand{\DecValTok}[1]{\textcolor[rgb]{0.25,0.63,0.44}{#1}}
\newcommand{\DocumentationTok}[1]{\textcolor[rgb]{0.73,0.13,0.13}{\textit{#1}}}
\newcommand{\ErrorTok}[1]{\textcolor[rgb]{1.00,0.00,0.00}{\textbf{#1}}}
\newcommand{\ExtensionTok}[1]{#1}
\newcommand{\FloatTok}[1]{\textcolor[rgb]{0.25,0.63,0.44}{#1}}
\newcommand{\FunctionTok}[1]{\textcolor[rgb]{0.02,0.16,0.49}{#1}}
\newcommand{\ImportTok}[1]{\textcolor[rgb]{0.00,0.50,0.00}{\textbf{#1}}}
\newcommand{\InformationTok}[1]{\textcolor[rgb]{0.38,0.63,0.69}{\textbf{\textit{#1}}}}
\newcommand{\KeywordTok}[1]{\textcolor[rgb]{0.00,0.44,0.13}{\textbf{#1}}}
\newcommand{\NormalTok}[1]{#1}
\newcommand{\OperatorTok}[1]{\textcolor[rgb]{0.40,0.40,0.40}{#1}}
\newcommand{\OtherTok}[1]{\textcolor[rgb]{0.00,0.44,0.13}{#1}}
\newcommand{\PreprocessorTok}[1]{\textcolor[rgb]{0.74,0.48,0.00}{#1}}
\newcommand{\RegionMarkerTok}[1]{#1}
\newcommand{\SpecialCharTok}[1]{\textcolor[rgb]{0.25,0.44,0.63}{#1}}
\newcommand{\SpecialStringTok}[1]{\textcolor[rgb]{0.73,0.40,0.53}{#1}}
\newcommand{\StringTok}[1]{\textcolor[rgb]{0.25,0.44,0.63}{#1}}
\newcommand{\VariableTok}[1]{\textcolor[rgb]{0.10,0.09,0.49}{#1}}
\newcommand{\VerbatimStringTok}[1]{\textcolor[rgb]{0.25,0.44,0.63}{#1}}
\newcommand{\WarningTok}[1]{\textcolor[rgb]{0.38,0.63,0.69}{\textbf{\textit{#1}}}}
\usepackage{longtable,booktabs,array}
\newcounter{none} % for unnumbered tables
\usepackage{calc} % for calculating minipage widths
% Correct order of tables after \paragraph or \subparagraph
\usepackage{etoolbox}
\makeatletter
\patchcmd\longtable{\par}{\if@noskipsec\mbox{}\fi\par}{}{}
\makeatother
% Allow footnotes in longtable head/foot
\IfFileExists{footnotehyper.sty}{\usepackage{footnotehyper}}{\usepackage{footnote}}
\makesavenoteenv{longtable}
\ifLuaTeX
\usepackage[bidi=basic,shorthands=off]{babel}
\else
\usepackage[bidi=default,shorthands=off]{babel}
\fi
\ifLuaTeX
  \usepackage{selnolig} % disable illegal ligatures
\fi
\setlength{\emergencystretch}{3em} % prevent overfull lines
\providecommand{\tightlist}{%
  \setlength{\itemsep}{0pt}\setlength{\parskip}{0pt}}
% Shared visual layer for the generated Kaleido FlowTwin reports.
\usepackage{helvet}
\renewcommand{\familydefault}{\sfdefault}
\usepackage{microtype}
\usepackage{xcolor}
\usepackage{titlesec}
\usepackage{fancyhdr}
\usepackage{enumitem}
\usepackage{booktabs}
\usepackage{xurl}
\usepackage{lastpage}
\usepackage{etoolbox}

\definecolor{KaleidoNavy}{HTML}{0A2342}
\definecolor{KaleidoBlue}{HTML}{176B87}
\definecolor{KaleidoCyan}{HTML}{3BB7C4}
\definecolor{KaleidoMint}{HTML}{DDF5F1}
\definecolor{KaleidoFog}{HTML}{F2F5F7}
\definecolor{KaleidoSlate}{HTML}{5D6A7A}
\definecolor{KaleidoCoral}{HTML}{F26B5B}

\titleformat{\section}
  {\Large\bfseries\color{KaleidoNavy}}
  {\thesection}{0.75em}{}[\vspace{0.2em}\color{KaleidoCyan}\titlerule]
\titleformat{\subsection}
  {\large\bfseries\color{KaleidoBlue}}
  {\thesubsection}{0.65em}{}
\titleformat{\subsubsection}
  {\normalsize\bfseries\color{KaleidoNavy}}
  {\thesubsubsection}{0.55em}{}
\titlespacing*{\section}{0pt}{2.2em}{1em}
\titlespacing*{\subsection}{0pt}{1.6em}{0.65em}

\pagestyle{fancy}
\fancyhf{}
\fancyhead[L]{\small\bfseries\color{KaleidoNavy}KALEIDO FLOWTWIN}
\fancyhead[R]{\small\color{KaleidoSlate}DOSSIER TECNICO}
\fancyfoot[L]{\scriptsize\color{KaleidoSlate}EVOCON Solutions GmbH}
\fancyfoot[R]{\scriptsize\color{KaleidoSlate}\thepage/\pageref*{LastPage}}
\renewcommand{\headrulewidth}{0.6pt}
\renewcommand{\headrule}{\hbox to\headwidth{\color{KaleidoCyan}\leaders\hrule height \headrulewidth\hfill}}
\renewcommand{\footrulewidth}{0pt}
\setlength{\headheight}{20pt}

\setlist[itemize]{leftmargin=1.4em,itemsep=0.25em,topsep=0.35em}
\setlist[enumerate]{leftmargin=1.7em,itemsep=0.25em,topsep=0.35em}
\setlength{\parindent}{0pt}
\setlength{\parskip}{0.55em}
\setlength{\emergencystretch}{3em}
\renewcommand{\arraystretch}{1.18}

\AtBeginEnvironment{quote}{%
  \color{KaleidoNavy}\itshape
  \setlength{\leftskip}{1em}
  \setlength{\rightskip}{1em}
}

\makeatletter
\renewcommand{\maketitle}{%
  \hypersetup{pageanchor=false}
  \begin{titlepage}
    \pagecolor{KaleidoNavy}
    \color{white}
    \vspace*{1.4cm}
    {\color{KaleidoCyan}\rule{2.4cm}{2.5pt}\par}
    \vspace{1.1cm}
    {\Huge\bfseries\@title\par}
    \vspace{0.65cm}
    {\Large\color{KaleidoMint}Kaleido FlowTwin\par}
    \vfill
    {\large Álvaro Schwiedop Souto\par}
    \vspace{0.15cm}
    {\normalsize\color{KaleidoMint}Industrial AI \& Predictive Quality Lead\par}
    \vspace{0.15cm}
    {\normalsize\color{KaleidoMint}EVOCON Solutions GmbH\par}
    \vspace{0.25cm}
    {\normalsize\color{KaleidoMint}\@date\par}
  \end{titlepage}
  \hypersetup{pageanchor=true}
  \nopagecolor
  \color{black}
}
\makeatother
\usepackage{bookmark}
\IfFileExists{xurl.sty}{\usepackage{xurl}}{} % add URL line breaks if available
\urlstyle{same}
% fallback for those not using the hyperref driver hyperxmp:
\makeatletter
\@ifundefined{xmpquote}{\newcommand{\xmpquote}[1]{#1}}{}
\makeatother
\hypersetup{
  pdftitle={SOTA 2026: world models{,} JEPA y procesos operativos},
  pdflang={es},
  colorlinks=true,
  linkcolor={KaleidoBlue},
  filecolor={Maroon},
  citecolor={Blue},
  urlcolor={KaleidoBlue},
  pdfcreator={LaTeX via pandoc}}

\title{SOTA 2026: world models, JEPA y procesos operativos}
\author{}
\date{15 julio 2026}

\begin{document}
\maketitle

{
\setcounter{tocdepth}{3}
\tableofcontents
}
\section{SOTA 2026: world models, JEPA y procesos
operativos}\label{sota-2026-world-models-jepa-y-procesos-operativos}

Fecha de corte: 15-07-2026. \texttt{SOTA} describe el estado de la
literatura, no el estado de FlowTwin. FlowTwin no tiene resultados.

\subsection{Tesis principal}\label{tesis-principal}

El avance 2026 no es un unico modelo ganador. Hay cinco tendencias:

\begin{enumerate}
\def\labelenumi{\arabic{enumi}.}
\tightlist
\item
  prediccion en estado latente en lugar de reconstruccion completa;
\item
  regularizacion explicita contra colapso;
\item
  acciones y horizontes multiples para convertir representacion en
  dinamica;
\item
  jerarquia, prefijos y planners para evitar rollouts largos costosos;
\item
  incertidumbre, fisica y adaptacion para despliegue bajo cambio.
\end{enumerate}

Para Kaleido, el SOTA de video/robotica aporta principios, no un modelo
que se pueda aplicar directamente. El dato es una secuencia de eventos
multiobjeto, irregular y parcialmente observada. Por ello el stack
inicial debe combinar process mining, prediccion de tiempo restante,
survival/boosting y un encoder de eventos. El modulo JEPA es una
hipotesis de valor incremental.

\subsection{Mapa de los 29 trabajos
indicados}\label{mapa-de-los-29-trabajos-indicados}

{\def\LTcaptype{none} % do not increment counter
\begin{longtable}[]{@{}lll@{}}
\toprule\noalign{}
Grupo & Trabajos & Leccion para FlowTwin \\
\midrule\noalign{}
\endhead
\bottomrule\noalign{}
\endlastfoot
Teoria SSL & Spectral SSL \texttt{2205.11508} & la relacion entre
vistas/eventos es supervision; una mala definicion aprende atajos \\
Features visuales & DINOv2 \texttt{2304.07193}, DINOv3
\texttt{2508.10104}, PatchCore \texttt{2106.08265} & usar backbones
congelados para fotos de incidencias; no entrenar vision desde cero \\
JEPA base & posicion H-JEPA \texttt{BZ5a1r-kVsf}, I-JEPA
\texttt{2301.08243} & predecir representaciones compatibles y separar
niveles temporales \\
Anticolapso & LeJEPA \texttt{2511.08544}, VISReg \texttt{2606.02572} &
medir varianza, rango, isotropia y probes; no confiar solo en loss \\
Incertidumbre & Var-JEPA \texttt{2603.20111} & futuros multimodales e
intervalos; comparar con conformal/ensembles simples \\
Temporal & TDV \texttt{2606.15956}, V-JEPA 2 \texttt{2506.09985}, V-JEPA
2.1 \texttt{2603.14482} & el tiempo y las acciones convierten features
en estado; features densas importan para localizacion \\
Fisica intuitiva & \texttt{2502.11831} & surprise puede detectar
violaciones, pero no prueba causalidad ni leyes explicitas \\
Implementacion & EB-JEPA \texttt{2602.03604} & componentes modulares y
small/single-GPU; util como referencia de tests \\
World model end-to-end & LeWorldModel \texttt{2603.19312},
identifiability \texttt{2605.26379} & una receta compacta es posible;
las garantias dependen de hipotesis que un proceso logistico no
satisface automaticamente \\
Rollouts & On Training in Imagination \texttt{2605.06732} & error de
dinamica y de coste se propaga; evaluar decisiones, no solo one-step
loss \\
Largo horizonte & HWM \texttt{2604.03208}, FF-JEPA \texttt{2606.09311},
Fast-LeWM \texttt{2606.26217} & jerarquia proyecto-turno-evento y
prediccion paralela por horizontes son mas utiles que rollout
autoregresivo profundo \\
Adaptacion & AdaJEPA \texttt{2606.32026} & recalibrar con transiciones
nuevas, pero conservar referencia congelada y trust region \\
Robotica & WorldDP \texttt{2606.08775}, SkyJEPA \texttt{2606.23444} &
separar plan de alto nivel y ejecucion; usar probes fisicos/operativos
interpretables \\
Generativo & MIRA, ahora arXiv \texttt{2607.05352} & la generacion
visual interactiva es cara y no necesaria para el MVP logistico \\
Multimodal & VL-JEPA \texttt{2512.10942}, MJEPA \texttt{2606.25225} &
texto/foto/audio pueden apoyar incidentes; cross-modal loss solo si hay
alineacion real \\
Ciencia & JEPA-DNA \texttt{2602.17162}, Phys-JEPA \texttt{2606.16076},
Neuro-JEPA \texttt{2606.14957} & objetivos hibridos, estado
fisico+residual, MoE y baselines simples; relevancia metodologica, no
transferencia directa \\
\end{longtable}
}

\subsection{Que aporta cada bloque}\label{que-aporta-cada-bloque}

\subsubsection{Representacion}\label{representacion}

I-JEPA muestra que ocultar bloques grandes y predecir embeddings puede
aprender semantica sin reconstruir pixeles. DINOv2/3 establece un
baseline visual mas practico para fotos de danos. En el MVP, la
fotografia es una modalidad opcional: se extraen features congeladas y
se prueba si anaden valor a eventos y notas. No se entrena un V-JEPA
sobre datos de Kaleido en fase 1.

\subsubsection{Dinamica}\label{dinamica}

Un encoder no es un world model. El salto requiere:

\begin{Shaded}
\begin{Highlighting}[]
\NormalTok{state\_t + verified\_action\_t + context\_t + horizon {-}\textgreater{} distribution(state\_future)}
\end{Highlighting}
\end{Shaded}

Las variables de receta, buque o tipo de carga son contexto. Asignar
equipo, cambiar secuencia, abrir un turno o reubicar un recurso pueden
ser acciones si estan timestamped y son intervenibles.

\subsubsection{Estabilidad}\label{estabilidad}

LeJEPA y LeWorldModel proponen regularizacion hacia una geometria
gaussiana. VISReg separa control de escala y forma. FlowTwin debe
registrar:

\begin{itemize}
\tightlist
\item
  desviacion estandar minima/media;
\item
  rango efectivo y participation ratio;
\item
  vecinos y separacion por proyecto;
\item
  probes congelados;
\item
  sensibilidad a shuffled actions/timestamps;
\item
  estabilidad por seed.
\end{itemize}

Una perdida predictiva baja puede coexistir con colapso o atajos.

\subsubsection{Horizonte}\label{horizonte}

Los proyectos logisticos son jerarquicos:

\begin{Shaded}
\begin{Highlighting}[]
\NormalTok{evento {-}\textgreater{} tarea {-}\textgreater{} operacion {-}\textgreater{} turno {-}\textgreater{} proyecto {-}\textgreater{} cadena}
\end{Highlighting}
\end{Shaded}

HWM/FF-JEPA/Fast-LeWM justifican modelar escalas y horizontes de forma
directa. La primera implementacion debe predecir 2/4/8 horas y fin de
operacion en paralelo. No necesita CEM ni planificar cientos de acciones
continuas.

\subsubsection{Adaptacion}\label{adaptacion}

AdaJEPA es atractiva para puertos, clientes y cargas diferentes. Sin
embargo, adaptar el unico modelo puede convertir retraso cronico en
normalidad. La arquitectura requerida es:

\begin{itemize}
\tightlist
\item
  referencia congelada;
\item
  adaptador/shadow por terminal o proceso;
\item
  gate de incertidumbre y distancia de update;
\item
  replay;
\item
  rollback/versionado;
\item
  reinicio explicito ante cambio de procedimiento.
\end{itemize}

\subsubsection{Incertidumbre}\label{incertidumbre}

Var-JEPA aporta una formulacion probabilistica, pero no es la unica
opcion. En un MVP data-poor deben compararse:

\begin{itemize}
\tightlist
\item
  quantile regression;
\item
  conformal prediction por grupos;
\item
  deep ensemble pequeno;
\item
  survival model con censura;
\item
  Var-JEPA solo si mejora calibracion/decision.
\end{itemize}

\subsection{SOTA mas cercano al dato de Kaleido: predictive process
monitoring}\label{sota-mas-cercano-al-dato-de-kaleido-predictive-process-monitoring}

El campo directamente aplicable es Predictive Process Monitoring (PPM):
predice siguiente actividad, siguiente timestamp, outcome y tiempo
restante desde prefijos de event logs.

Referencias de 2024-2026:

\begin{itemize}
\tightlist
\item
  PGTNet transforma logs a grafos y predice remaining time sobre 20
  logs: \url{https://arxiv.org/abs/2404.06267}.
\item
  La comparacion object-centric vs clasica resalta el valor de OCEL en
  procesos logisticos multiobjeto:
  \url{https://link.springer.com/article/10.1007/s10115-025-02461-y}.
\item
  Object-centric graph embeddings predicen next activity/time:
  \url{https://arxiv.org/abs/2507.15411}.
\item
  Un foundation model in-context para logs heterogeneos propone
  adaptacion con un pequeno support set, relevante para cold start:
  \url{https://doi.org/10.1109/ACCESS.2026.3658877}.
\item
  Un caso real de almacen outbound con 169.523 traces encuentra que deep
  models ganan en precision, pero boosting es competitivo con menor
  coste: \url{https://arxiv.org/abs/2509.18986}.
\item
  La version retrieval-augmented publicada como preprint el 09-07-2026
  es prometedora pero no peer-reviewed:
  \url{https://www.preprints.org/manuscript/202607.0705}.
\end{itemize}

Decision: implementar primero baselines de PPM y object-centric graph.
Evaluar un encoder JEPA de eventos despues. No llamar foundation model a
un Transformer pequeno entrenado sobre un solo log.

\subsection{Arquitectura SOTA-target, no SOTA
claim}\label{arquitectura-sota-target-no-sota-claim}

\begin{Shaded}
\begin{Highlighting}[]
\NormalTok{OCEL{-}like event graph}
\NormalTok{  {-}\textgreater{} categorical/time/value encoders}
\NormalTok{  {-}\textgreater{} object{-}centric temporal graph state}
\NormalTok{  {-}\textgreater{} multi{-}horizon heads: remaining time, deviation, next event}
\NormalTok{  {-}\textgreater{} optional JEPA target: future operation{-}state embedding}
\NormalTok{  {-}\textgreater{} calibrated uncertainty}
\NormalTok{  {-}\textgreater{} advisory scenario ranking over approved actions}
\end{Highlighting}
\end{Shaded}

Loss candidata:

\begin{Shaded}
\begin{Highlighting}[]
\NormalTok{L = L\_remaining\_quantile}
\NormalTok{  + lambda\_risk * L\_risk}
\NormalTok{  + lambda\_next * L\_next\_event}
\NormalTok{  + lambda\_jepa * L\_future\_embedding}
\NormalTok{  + lambda\_cal * L\_calibration}
\NormalTok{  + lambda\_reg * L\_anticollapse}
\end{Highlighting}
\end{Shaded}

El termino JEPA se activa solo tras demostrar señal y volumen. Cada
termino se abla por separado.

\subsection{Lo que no demuestra la
literatura}\label{lo-que-no-demuestra-la-literatura}

\begin{itemize}
\tightlist
\item
  Que un JEPA mejore automaticamente datos logisticos tabulares.
\item
  Que next-event accuracy implique ahorro.
\item
  Que una prediccion observacional sea causal.
\item
  Que video generado sea un simulador de operacion fiel.
\item
  Que adaptacion online sea segura sin referencia congelada.
\item
  Que un resultado de robotica transfiera a puertos.
\item
  Que mas parametros ganen con pocos proyectos.
\end{itemize}

\subsection{Gate para conservar JEPA}\label{gate-para-conservar-jepa}

Se conserva el modulo si, con el mismo split y presupuesto:

\begin{enumerate}
\def\labelenumi{\arabic{enumi}.}
\tightlist
\item
  mejora el error de tiempo restante o AUPRC de desviacion frente a
  boosting y ProcessTransformer;
\item
  mejora al menos un KPI operativo: lead time o falsas alarmas;
\item
  la mejora aparece en tres seeds y en holdout por proyecto/cliente;
\item
  correct actions superan shuffled actions cuando se hable de accion;
\item
  no empeora calibracion materialmente;
\item
  el coste de inferencia cabe en el SLA.
\end{enumerate}

Si no cumple, el producto se entrega con el modelo simple ganador.

\subsection{Referencias primarias indicadas por el
usuario}\label{referencias-primarias-indicadas-por-el-usuario}

La bibliografia completa esta en \texttt{REFERENCES.bib}. Los enlaces de
entrada se conservan en \texttt{Apuntes\_IA/Referencias\ ArXiv.txt}. Los
apuntes locales \texttt{mega\_apuntes\_jepa\_world\_models.tex} ya
organizan los 29 trabajos en orden conceptual y han sido usados como
mapa, contrastando titulos/resumenes con arXiv a fecha de corte.

\end{document}
